\documentclass[12pt]{article}

\usepackage[T2A]{fontenc}
\usepackage[utf8]{inputenc}
\usepackage[russian]{babel}
\usepackage{graphicx}
\usepackage{amsmath}
\usepackage{amssymb}
\usepackage{float}
\usepackage{booktabs}
\usepackage{geometry}

\geometry{
    left=30mm,
    right=15mm,
    top=20mm,
    bottom=20mm
}

\title{Исследование связи индивидуальной волатильности и взаимной корреляции компонентов с волатильностью индекса S\&P 500}
\author{Леонид Кононенко}
\date{Сентябрь 2026}

\begin{document}

\maketitle

\begin{abstract}

В работе исследуется связь между волатильностью индекса S\&P 500 и характеристиками доходностей акций, входящих в его состав. Основное внимание уделяется двум факторам: индивидуальной волатильности компонентов индекса и взаимной корреляции между их доходностями. Исходные данные представлены историческими дневными котировками акций компонентов S\&P 500 и самого индекса.

В ходе исследования выполняется сбор и предварительная обработка финансовых данных, рассчитываются логарифмические доходности, скользящие оценки волатильности и средняя попарная корреляция между компонентами индекса. Подготовленные данные сохраняются в реляционной базе данных SQLite, после чего проводится их статистический и визуальный анализ.

Для количественной оценки взаимосвязей строятся модели множественной линейной регрессии, в которых зависимой переменной выступает волатильность S\&P 500, а в качестве объясняющих переменных используются средняя волатильность компонентов, средняя взаимная корреляция, абсолютная доходность индекса и средний объём торгов. Качество моделей оценивается с помощью коэффициента детерминации $R^2$, скорректированного коэффициента детерминации, MAE и RMSE.

Целью исследования является определение того, насколько индивидуальная волатильность компонентов и их взаимная корреляция связаны с изменением общей волатильности индекса S\&P 500.

\end{abstract}

\newpage

\tableofcontents

\newpage

\section{Введение}

Фондовый рынок характеризуется постоянным изменением цен финансовых активов и неопределённостью относительно их будущей динамики. Одним из основных показателей, используемых для количественной оценки уровня неопределённости и риска на финансовом рынке, является волатильность. Анализ волатильности имеет важное значение при оценке риска, формировании инвестиционных портфелей и исследовании поведения финансовых рынков.

Особый интерес представляет изучение волатильности широких фондовых индексов. В отличие от отдельной акции, индекс отражает динамику большого количества компаний, относящихся к различным секторам экономики. Одним из наиболее известных фондовых индексов является S\&P 500, который используется как один из основных индикаторов состояния американского фондового рынка.

Волатильность индекса определяется не только тем, насколько сильно изменяются цены отдельных компаний, но и тем, насколько синхронно эти изменения происходят. Если доходности большого количества акций начинают двигаться в одном направлении, взаимная корреляция между ними увеличивается. В результате даже при неизменной индивидуальной волатильности отдельных компонентов общая волатильность индекса может возрастать.

Таким образом, возникает вопрос о том, какую роль играют индивидуальная волатильность компонентов и взаимная корреляция между ними в формировании волатильности S\&P 500. Данный вопрос представляет интерес с точки зрения статистического анализа финансовых данных и позволяет связать характеристики отдельных акций с поведением агрегированного рыночного индикатора.

Актуальность исследования обусловлена широким использованием фондовых индексов для оценки состояния финансового рынка и необходимостью количественного анализа факторов, связанных с изменением рыночного риска. Кроме того, современные источники финансовых данных позволяют работать с большими массивами исторических котировок, что создаёт возможность для практического применения методов статистического анализа и регрессионного моделирования.

Объектом исследования является индекс S\&P 500 и акции компаний, входящих в его состав.

Предметом исследования являются статистические взаимосвязи между волатильностью индекса S\&P 500, индивидуальной волатильностью его компонентов и взаимной корреляцией между их доходностями.

Целью работы является исследование связи между индивидуальной волатильностью компонентов, их взаимной корреляцией и волатильностью индекса S\&P 500.

Для достижения поставленной цели необходимо решить следующие задачи:

\begin{enumerate}
    \item собрать исторические данные о ценах и объёмах торгов акций компонентов S\&P 500;
    \item собрать исторические данные по индексу S\&P 500;
    \item выполнить очистку и предварительную обработку исходных данных;
    \item организовать подготовленные данные в реляционной базе данных SQLite;
    \item рассчитать логарифмические доходности отдельных акций и индекса;
    \item рассчитать скользящие показатели волатильности;
    \item определить среднюю волатильность компонентов индекса;
    \item рассчитать среднюю попарную корреляцию между доходностями компонентов;
    \item провести первичный статистический анализ полученных данных;
    \item выполнить визуальный анализ динамики волатильности и корреляций;
    \item построить модели множественной линейной регрессии;
    \item оценить качество построенных моделей;
    \item сравнить вклад индивидуальной волатильности и взаимной корреляции в объяснение волатильности S\&P 500.
\end{enumerate}

В качестве источника исторических финансовых данных используется Yahoo Finance. Сбор и обработка данных осуществляются с помощью языка программирования Python. Для хранения подготовленных данных используется реляционная база данных SQLite.

В качестве основной характеристики изменения цены используется логарифмическая доходность:

\begin{equation}
r_{i,t} =
\ln\left(\frac{P_{i,t}}{P_{i,t-1}}\right),
\end{equation}

где $P_{i,t}$ --- цена актива $i$ в момент времени $t$.

Для оценки волатильности используется скользящее стандартное отклонение доходностей. В частности, при использовании окна длиной 20 торговых дней волатильность отдельной акции может быть рассчитана следующим образом:

\begin{equation}
\sigma_{i,t}
=
\sqrt{252}
\sqrt{
\frac{1}{19}
\sum_{k=0}^{19}
\left(r_{i,t-k}-\bar r_{i,t}\right)^2
},
\end{equation}

где $\bar r_{i,t}$ --- средняя доходность за рассматриваемое 20-дневное окно, а множитель $\sqrt{252}$ используется для перевода дневной волатильности в годовой масштаб.

Для характеристики общего уровня волатильности компонентов вводится показатель средней волатильности:

\begin{equation}
MeanVol_t =
\frac{1}{N}
\sum_{i=1}^{N}\sigma_{i,t},
\end{equation}

где $N$ --- количество рассматриваемых компонентов индекса.

Для характеристики степени синхронности движения акций используется средняя попарная корреляция:

\begin{equation}
\overline{\rho}_t =
\frac{2}{N(N-1)}
\sum_{i<j}\rho_{ij,t},
\end{equation}

где $\rho_{ij,t}$ --- коэффициент корреляции между доходностями акций $i$ и $j$ за рассматриваемый временной интервал.

В рамках регрессионного анализа волатильность индекса S\&P 500 рассматривается в качестве зависимой переменной. В качестве факторов используются средняя волатильность компонентов, средняя взаимная корреляция, абсолютная доходность индекса и средний объём торгов. Общий вид модели имеет следующий вид:

\begin{equation}
\sigma_{S\&P500,t}
=
\beta_0
+
\beta_1 MeanVol_t
+
\beta_2 \overline{\rho}_t
+
\beta_3 |r_{S\&P500,t}|
+
\beta_4 MeanVolume_t
+
\varepsilon_t,
\end{equation}

где $\beta_0$ --- свободный член, $\beta_1,\ldots,\beta_4$ --- коэффициенты регрессии, а $\varepsilon_t$ --- случайная ошибка.

Полученные результаты позволяют оценить статистическую связь между характеристиками отдельных компонентов индекса и его общей волатильностью. При этом регрессионный анализ в рамках данной работы рассматривается как способ выявления статистических взаимосвязей, а не как доказательство причинно-следственных отношений.

Работа состоит из трёх основных разделов. В первом разделе рассматриваются исходные данные, методы их сбора и обработки, структура базы данных и результаты первичного статистического анализа. Во втором разделе проводится визуальный анализ исследуемых показателей. Третий раздел посвящён построению и оценке регрессионных моделей волатильности индекса S\&P 500.

\newpage

\section{Сбор и первичный анализ данных}

\subsection{Постановка задачи и характеристика исследуемого объекта}

\subsection{Источники и структура исходных данных}

\subsection{Методика сбора данных}

\subsection{Проектирование базы данных}

\subsection{Предварительная обработка данных}

\subsection{Расчёт доходностей и показателей волатильности}

\subsection{Первичный статистический анализ данных}

\newpage

\section{Визуальный анализ данных}

\subsection{Динамика индекса S\&P 500 и его волатильности}

\subsection{Сравнение волатильности компонентов индекса}

\subsection{Анализ волатильности по секторам экономики}

\subsection{Корреляционный анализ доходностей компонентов}

\subsection{Интерпретация результатов визуального анализа}

\newpage

\section{Регрессионное моделирование волатильности S\&P 500}

\subsection{Постановка регрессионной задачи}

\subsection{Формирование зависимой и независимых переменных}

\subsection{Разделение данных на обучающую и тестовую выборки}

\subsection{Построение моделей множественной линейной регрессии}

\subsection{Оценка качества моделей}

\subsection{Интерпретация коэффициентов регрессии}

\subsection{Сравнение моделей}

\newpage

\section*{Заключение}
\addcontentsline{toc}{section}{Заключение}

\newpage

\section*{Список использованных источников}
\addcontentsline{toc}{section}{Список использованных источников}

\newpage

\section*{Приложения}
\addcontentsline{toc}{section}{Приложения}

\end{document}